Since @@M0@@, we have @@M1@@, and consequently:
@@M2@@
Put:
@@M3@@
Since @@M4@@, we have:
@@M5@@
By V (6.2.1), this gives:
@@M6@@
Taking 3.5 into account, we finally obtain:
@@M7@@

By the definition of @@M8@@, we therefore have:
\sgaStableTarget{sga6.e08.eq.5-3-1}@@M9@@

\sgaStableTarget{sga6.e08.para.5-4}\textbf{5.4.} Using the basis of @@M10@@ over @@M11@@ formed by the elements @@M12@@, with @@M13@@, one sees by additivity that it suffices to prove the theorem for elements of the form @@M14@@. We then have:
@@M15@@
On the other hand:
@@M16@@

One thus sees that it suffices to prove the theorem for the @@M17@@, with @@M18@@. Now by (VI (5.2.3)), we have
@@M19@@
Consequently:
\sgaStableTarget{sga6.e08.eq.5-4-1}@@M20@@
